\section{Training with a measured degradation}
\label{sec:method}
We keep a published real-time backbone and change its fine-tuning in two respects: the degradation applied to the training images (Section~\ref{sec:measured}) and the photometric range of the training data (Section~\ref{sec:stable}).

\subsection{Degradation of real small ear images}
\label{sec:measured}
We use the images of the \RoleFit{} EarVN1.0 subjects reserved for this purpose, which comprise \FitSmall{} real small images (short side 24 to 48 pixels) and \FitLarge{} large images. Three quantities are estimated.

\paragraph{Compression level} A JPEG file stores its quantisation tables, from which the quality setting of a standard encoder can be recovered. In the real small images the distribution is bimodal: \QLowShare{}\% of the images are at quality 75 and \QHighShare{}\% at quality 93, with the remainder at neighbouring values. We use this empirical distribution.

\paragraph{Noise} We apply the estimator of Immerk{\ae}r~\cite{immerkaer1996noise} to each real small image and use the 5th to 95th percentile range, \NoiseLo{} to \NoiseHi{} on the 8-bit scale. Because the images are already compressed, this range underestimates the noise before compression.

\paragraph{Blur} A blur kernel is difficult to identify from a compressed image of a few dozen pixels. We select, from a set of candidate ranges for the width of an isotropic Gaussian, the range for which the sharpness statistics (Laplacian energy normalised by variance) of large images degraded with the full pipeline are closest to those of the real small images in Kolmogorov--Smirnov distance. The selected range is $\sigma\in[\BlurLo,\BlurHi]$ input pixels, which corresponds to a small amount of blur.

The resulting pipeline consists of Gaussian blur, bicubic downsampling, Gaussian noise and JPEG compression, with parameters drawn independently for each training sample. The generic pipeline of Section~\ref{sec:degradations} has the same structure with wider parameter ranges.

\paragraph{Comparison with real images} To assess how close the simulated images are to real ones, we train a small convolutional classifier to separate $16\times16$ patches of real small ear images from patches of simulated images, which are generated from large images of the same subjects and brought to the same size distribution. Subjects are split between training and testing (\RealTrainSubj{} and \RealTestSubj{} subjects of the classifier group), and the experiment is repeated with three seeds. Table~\ref{tab:realism} gives the balanced accuracy, where 50\% means that the classifier does not separate simulated from real images. The classifier separates bicubic downsampling (\RealBic{}\%) and the generic pipeline (\RealGen{}\%) from real images with similar accuracy, so by this test the generic pipeline is not closer to the real data than bicubic downsampling. The measured pipeline gives \RealEst{}\%. Bicubic downsampling followed by JPEG at quality 75 gives \RealJpg{}\%, which is not significantly different from the measured pipeline. By this test, compression is the main factor that makes the simulated images similar to real ones. We return to this point in Section~\ref{sec:ablation}.

\begin{table}[t]
\centering
\caption{Balanced accuracy (\%) of a classifier trained to separate real small ear images from simulated ones, on held-out subjects. Closer to 50\% means more realistic. Intervals are 95\% bootstrap intervals.}
\label{tab:realism}
\small
\begin{tabular}{@{}lcc@{}}
\toprule
Simulated with & Accuracy & 95\% CI \\
\midrule
\tabinput{generated/tab_realism}
\bottomrule
\end{tabular}
\end{table}

\subsection{Stability of fine-tuning}
\label{sec:stable}
Our first models were fine-tuned on AMI alone. Their validation PSNR on AMI improved, but on EarVN1.0 the SPAN model fine-tuned with bicubic downsampling reached \CollapsePsnr{}~dB on clean inputs, compared with \CollapsePub{}~dB for the published weights it was initialised from and \CollapseBic{}~dB for bicubic interpolation. \CollapseBroken{} of the \NEarvnImg{} outputs were below 20~dB and contained high-amplitude noise. The failures occurred mainly on images with bright regions. AMI was acquired under uniform indoor lighting and contains few such images. Our inspection of the network suggests that, after fine-tuning, the activations of the attention blocks grow from block to block for inputs brighter than the training images. Validation on held-out AMI subjects did not reveal the problem.

Two changes prevented this failure in our experiments. The first is a photometric augmentation: with probability 0.75 the ground-truth patch is multiplied by a random gain (log-uniform in $[0.7, 2.2]$), raised to a random gamma ($[0.7,1.4]$) and scaled per channel ($[0.9,1.1]$) before the degradation is applied. The second is the addition of in-the-wild training images: 30\% of the training samples are drawn from the large images of the EarVN1.0 training subjects. We also add a stress test to each training run. The validation images are brightened by factors of 1.4 and 1.8, and a run is marked as unstable if any output is more than 10~dB below bicubic interpolation or if the margin of the model over bicubic interpolation decreases by more than 1.5~dB. All \NRuns{} runs reported in this paper pass this test (largest decrease \StressMaxDrop{}~dB). Table~\ref{tab:stability} shows the effect of the two changes.

\begin{table}[t]
\centering
\caption{Effect of the fine-tuning recipe on clean inputs (PSNR-Y, dB): SPAN fine-tuned with bicubic downsampling. $^\dagger$One run, excluded from all other results.}
\label{tab:stability}
\small
\begin{tabular}{@{}lcccc@{}}
\toprule
 & AMI & EarVN1.0 & AWEx & AWEx \\
 & 36 px & 24 px & 24 px & 36 px \\
\midrule
\tabinput{generated/tab_stability}
\bottomrule
\end{tabular}
\end{table}

\subsection{Training details}
We fine-tune from the published weights for 20{,}000 iterations with the $L_1$ loss, Adam~\cite{kingma2015adam} with an initial learning rate of $2\times10^{-4}$ and cosine decay, batches of 64 input patches of $24\times24$ pixels, and an exponential moving average of the weights (decay 0.999). Each AMI training image is first resized so that its short side is drawn at random between 96 and 192 pixels, so that one model covers the three benchmark sizes. The checkpoint with the highest PSNR on the validation subjects of the fold is kept. The random seed of a run equals its fold index. One run takes about half an hour on an RTX~3080. We apply the procedure to two backbones from different architecture families: SPAN with 48 channels~\cite{wan2024span} and DISP, an entry of the NTIRE 2026 challenge~\cite{ren2026ntire}.
